We began with a question that sounded technical: \emph{can a small
language model learn to solve programming tasks?} The answer we got was
more interesting and less about the model than we expected.

\section{Where the journey started}\label{where-the-journey-started}

The premise was: building a small model is relatively easy; determining
what it actually learned is the hard part. That premise turned out to be
more accurate than we had any right to expect. The model was built in an
afternoon. Understanding what it had --- and had not --- learned took
the rest of the book.

That first line deserves to be unpacked, because it is the whole book in
miniature. A decoder-only transformer that maps a handful of programming
sentences to Python is an afternoon's work once the pieces are in place.
But \emph{knowing what that model actually did} --- whether it learned,
memorized, associated, or cheated --- required a world, a teacher, a
judge, rivals, and a battery of controlled experiments. The model was
the easy part. The certainty about what it did was the hard part, and it
is the part worth writing about.

We built a laboratory, not just a model. Four artifacts: a
\textbf{controlled world} (an AST-whitelisted Python sandbox), a
\textbf{teacher} (a semantic dataset generator), a \textbf{judge} (an
execution-based, mutation-tested evaluator), and the \textbf{student}
(the model itself). The model was only one component. The experiment was
the subject.

\section{What the laboratory
produced}\label{what-the-laboratory-produced}

The experiments told a clean, coherent, honest story:

\begin{itemize}
\tightlist
\item
  \textbf{The model learns.} On individual operations it reliably maps
  natural language to correct Python, beating a hand-written template
  rule system and reaching \textasciitilde90\% in-distribution.
\item
  \textbf{But it is surface-tied.} Change the variable names, and it
  drops to 0.35. Change the wording while keeping the computation
  identical, and it falls to \textbf{0.000} --- below both baselines.
\item
  \textbf{And it cannot compose.} Even though it knows every primitive,
  it scores \textbf{0.000} on operations that combine two things it
  knows --- at every model size we tested, up to 29.3M.
\item
  \textbf{Neither capacity nor data fixes it at this scale.} A
  \textasciitilde33× larger model (up to 29.3M) and many more examples
  leave the composition boundary at 0.000.
\end{itemize}

Along the way, the laboratory caught itself. The first time we evaluated
the model, it scored 0.000 --- and it wasn't the model's fault. The
evaluator was broken (a NUL- padding artifact in the output made every
program fail to parse). Fixing the \emph{ tester} revealed the model's
true 0.897. That moment --- discovering our own instrument was wrong
before we trusted it --- is as much a part of this book as any model
result.

\section{The deeper achievement}\label{the-deeper-achievement}

The book's real finding is not ``small models are limited.'' It is:

\begin{quote}
\textbf{Building a trustworthy experiment was harder and more valuable
than building the model.}
\end{quote}

The model was the instrument. The experiment was the subject. The things
that made this work --- the sandbox that rejects anything not explicitly
allowed, the judge that executes code and compares semantics rather than
strings, the mutation suite that red-teams the judge, the leakage audit
that caught our own dataset leaking the answer, the baselines that ask
``could a simpler machine do this?'', the four-level taxonomy that
refuses to collapse everything into one number --- those are the actual
achievements here.

It is worth lingering on why each of these matters, because together
they form a discipline rather than a checklist. The sandbox's allowlist
(not a blocklist) meant the model could never do something we did not
intend, so a failure could always be attributed to the model rather than
to an accident of environment. The judge's execution-based, semantic
comparison meant a ``correct'' answer was objectively correct, not
merely stylistically close. The mutation suite meant we doubted the
judge before we trusted it. The baseline comparisons meant a number only
had meaning relative to a simpler machine. And the four-level taxonomy
meant we never overstated a result by collapsing different kinds of
generalization into one score.

These are the disciplines that make a claim about ``what a model
learned'' worth trusting. They are the real product of this book.

\section{The recurring habit}\label{the-recurring-habit}

There is a way of thinking that runs through everything here, and it is
worth naming directly because it is what a careful investigation
actually looks like. It is the habit of distrusting the easy answer ---
including our own.

Every time we got a number, the next step was not to celebrate or to
move on. It was to ask: \textbf{how could the model have gotten this
without learning what we think it learned?} Could it have memorized?
Retrieved? Exploited a template? Leaked the answer from our own data?
Hit a bug in our evaluator? Each of those is a way a
seemingly-impressive result can be hollow, and each one we then built a
control to test. That constant, deliberate suspicion --- aimed at our
own apparatus as much as at the model --- is the difference between ``we
ran an experiment'' and ``we know what the experiment means.''

\section{What the reader takes away}\label{what-the-reader-takes-away}

If we take one thing from this book, it is a habit, not a fact:

\begin{quote}
\textbf{Never trust a number from a model we don't understand, or from
an experiment we haven't audited.}
\end{quote}

Before we believe ``the model learned to add,'' we have to be able to
answer: remove the memorization, the retrieval, the templates, the
leakage, the evaluator artifacts --- what remains? In our case, for
simple primitives, what remained was genuine
learning-within-a-surface-form. It was real, and it was narrow, and it
did not compose.

That is the honest, complete answer. We did not build a general
programmer. We built the means to ask, precisely and honestly, what a
small model can and cannot learn --- and we got an answer worth having.

\textbf{Build the model. Test the model. Break the evaluator. Break the
model. Audit the data. Challenge the conclusion. Repeat.} Only then say:
``This is what the model learned.''

The model was not the point. The experiment was. Because the world is
full of models, but a disciplined way of finding out what they actually
learned --- that is rare, and that is what is worth keeping.
